\subsection{29. Divide and Conquer: Learning Chaotic Dynamical Systems with Multi-Step Penalty Neural ODEs (MP-NODE)}
\textbf{OSTI:} \texttt{3003857}\quad\textbf{Authors:} Chakraborty, Chung, Arcomano, Maulik (2024)\quad\textbf{Domain:} Scientific ML / Chaotic Dynamical Systems\\
\scorebadge{Coverage}{5}\quad\scorebadge{Agreement}{4}\quad\textbf{Total:} 9/20\par\smallskip
\noindent\textbf{Coverage rationale.} Reproduced Section 4.1.1 (gradient explosion), 4.1.2 (loss landscape, partial), Lorenz-63 vanilla vs MP-NODE, and attempted KS equation. NOT replicated: Kolmogorov flow (Section 4.3), ERA5 climate emulation (Section 4.4), which are the paper's headline high-dimensional results. Used PyTorch + RK4 vs paper's likely JAX + Tsit5; architecture choices (Lorenz MLP, KS dilated-CNN) partially reconstructed.\par
\noindent\textbf{Agreement rationale.} Gradient reduction qualitatively reproduced but quantitatively much smaller (\textasciitilde{}40x vs paper's \textasciitilde{}10\textasciicircum{}8x) because T=10 vs paper's T=20 is too short for full chaotic amplification. MP-NODE beats vanilla on Lorenz training loss (22.5 vs 81.2). Attractor statistics mismatch: our MP mean z=97 vs true 23.5 (neither model fully converged). KS experiment collapsed to zero output --- effectively failed. So MP \textgreater{} vanilla qualitative claim holds; quantitative claims only partially.\par\smallskip
\noindent\textbf{What reproduced:}
\begin{itemize}[leftmargin=1.4em,itemsep=0.2em,topsep=0.2em]
\item Multi-step penalty augmented loss formulation (Algorithm 1, Eqs. 16-24)
\item Gradient-magnitude reduction of MP vs vanilla (qualitative)
\item Similar computational cost (MP \textasciitilde{}5\% overhead)
\item MP-NODE lower L\_GT vs vanilla on Lorenz-63 (3.6x lower)
\item Multi-step window + mu-scheduling machinery
\item Vanilla NODE attractor collapse vs MP-NODE retaining orbital structure
\end{itemize}\noindent\textbf{What missing:}
\begin{itemize}[leftmargin=1.4em,itemsep=0.2em,topsep=0.2em]
\item Section 4.3 Kolmogorov flow (2D NS turbulence)
\item Section 4.4 ERA5 atmospheric emulation
\item Quantitative gradient explosion 10\textasciicircum{}8 factor (hardware-limited)
\item Working KS joint-PDF reproduction (architecture collapse)
\item Paper Table 1 KL divergence 0.029 for MP-NODE on KS
\item Original JAX implementation; stabilized CNN-NODE of Linot 2023
\item Rigorous loss-landscape non-convexity demonstration
\end{itemize}\noindent\textbf{Follow-on research questions:}
\begin{enumerate}[leftmargin=1.6em,itemsep=0.2em,topsep=0.2em,label=Q\arabic*.]
\item Rerun the gradient explosion demo at T=20 and 2000 steps (matching paper) on a single A100 / Spark to confirm the 10\textasciicircum{}8x gap --- does MP's advantage scale with T as exp(lambda\_1 * T)
\item Implement Linot-2023's stabilized NODE architecture with dilated convolutions (rates 1,2,3,4,3,2,1) and retest the KS joint-PDF; does this rescue the KS training collapse seen in the current replication
\item Compare MP-NODE against LSS (least-squares shadowing) on Lorenz-63 for exact gradient accuracy at fixed compute --- is MP's O(m+n) cost vs LSS's O(m*n\textasciicircum{}3) the only justification or is MP also more accurate
\item Study the mu schedule automatically: can we derive an adaptive mu update rule (e.g. based on the penalty/GT loss ratio) that eliminates the 'problem-dependent empirical' schedule the paper describes
\item Apply MP-NODE to a shallow-water / Lorenz-96 / Kuramoto system and test whether a single mu schedule transfers across problems or truly requires per-system tuning.
\end{enumerate}

